% ============================================================================
% Fichier  : partie4/P04_C12_forensique_systeme.tex
% Titre    : Forensique Système
% Auteur   : MINKA MI NGUIDJOÏ Thierry Emmanuel
% ============================================================================
% Sources : NIST SP 800-86 ; artefacts Windows et Linux. Opération FARO.
% ============================================================================

\chapter{Forensique Système}
\label{chap:systeme}
\chaptermark{Forensique système}

\epigraph{%
    « Le disque garde la mémoire de ce que l'utilisateur croyait avoir effacé. »%
}{Principe du cours}

\niveauChapitre{Intermédiaire à expert}{Chapitre~\ref{chap:fondements}~;
administration Windows et Linux.}

\begin{boxobjectifs}
\begin{itemize}
    \item Identifier les artefacts système utiles sous Windows et Linux.
    \item Reconstituer une chronologie à partir des horodatages.
    \item Repérer les traces de persistance et d'exécution.
    \item Distinguer suppression logique et effacement réel.
    \item Rédiger un constat système.
\end{itemize}
\end{boxobjectifs}

\bigskip

% ============================================================================
\section*{Question directrice : que révèle le disque de FS-01 sur les onze jours~?}
% ============================================================================

Le disque de \texttt{FS-01} contient les traces du passage de l'attaquant~:
comptes créés, services installés, programmes exécutés, fichiers consultés.
Ce chapitre apprend à les lire et à en tirer une chronologie opposable.

\medskip
\textbf{Plan du chapitre.}
\begin{enumerate}
    \item Les artefacts Windows (section~\ref{sec:c12-windows}).
    \item Les artefacts Linux (section~\ref{sec:c12-linux}).
    \item La chronologie (section~\ref{sec:c12-chrono}).
    \item Suppression et récupération (section~\ref{sec:c12-suppression}).
\end{enumerate}

\separateur

% ============================================================================
\section{Les Artefacts Windows}
\label{sec:c12-windows}
% ============================================================================

\begin{tableaucours}{L{4.0cm}L{8.2cm}}{%
    \tableauHeader{Artefact} &
    \tableauHeader{Ce qu'il révèle}
}
Journaux d'événements & Connexions, comptes, services, tâches \\
Base de registre & Programmes au démarrage, périphériques branchés \\
Traces d'exécution & Programmes lancés, fréquence, dates \\
Fichiers récents, raccourcis & Documents ouverts, chemins \\
Fichier d'échange, clichés & Fragments de mémoire, versions antérieures \\
\end{tableaucours}
\vspace{-0.8em}
\begin{center}{\small\itshape Tableau~12.1. Artefacts forensiques Windows}\end{center}

% ============================================================================
\section{Les Artefacts Linux}
\label{sec:c12-linux}
% ============================================================================

\begin{tableaucours}{L{4.0cm}L{8.2cm}}{%
    \tableauHeader{Artefact} &
    \tableauHeader{Ce qu'il révèle}
}
Journaux système & Authentifications, sudo, démons \\
Historique des interpréteurs & Commandes saisies \\
Tâches planifiées (cron) & Persistance, exécutions régulières \\
Horodatages des fichiers & Accès, modification, changement de métadonnées \\
Comptes et clés & Comptes ajoutés, clés d'accès déposées \\
\end{tableaucours}
\vspace{-0.8em}
\begin{center}{\small\itshape Tableau~12.2. Artefacts forensiques Linux}\end{center}

% ============================================================================
\section{La Chronologie}
\label{sec:c12-chrono}
% ============================================================================

\begin{boxprocedure}
\textbf{Construire une chronologie (\textit{timeline})}
\begin{enumerate}
    \item Rassembler les horodatages de tous les artefacts.
    \item Recaler les horloges (écarts connus).
    \item Ordonner les événements et repérer les enchaînements.
    \item Distinguer le certain (trace) du probable (interprétation).
    \item Dater chaque étape de l'attaque.
\end{enumerate}
\end{boxprocedure}

\begin{boxscenario}{Opération FARO. Onze jours reconstitués}
\textbf{Constatation.} La chronologie du disque de \texttt{FS-01} établit~:
23~février, connexion du compte \texttt{svc\_sauvegarde} à 2~h~14~; 24~février,
création du compte \texttt{svc\_maj2}~; 1\textsuperscript{er}~mars,
installation du service \texttt{WinUpd32}~; 4~mars, accès massif aux
partages~; 6~mars, chiffrement.

\medskip
\textbf{Analyse} (confiance élevée pour les traces, moyenne pour les
intentions). Les horodatages, recalés, dessinent une progression cohérente.
On distingue ce qui est tracé (les événements) de ce qu'on en déduit (le but
de l'attaquant).
\end{boxscenario}

% ============================================================================
\section{Suppression et Récupération}
\label{sec:c12-suppression}
% ============================================================================

Supprimer un fichier ne l'efface pas~: le système marque l'espace comme
réutilisable, mais les données subsistent jusqu'à réécriture. D'où la
possibilité de \termecle{récupérer} des fichiers supprimés, tant qu'ils n'ont
pas été écrasés. L'attaquant qui « nettoie » laisse souvent des traces
récupérables.

\begin{boxmetier}{Analyste forensique}
\textbf{Séparer le tracé du déduit.} Dans un constat système, l'analyste
écrit d'abord ce que la trace montre, puis, à part, ce qu'il en déduit, avec
un niveau de confiance. C'est le cœur de la méthode constatation/analyse.
\end{boxmetier}

\separateur

\begin{boxresume}
\textbf{Ce qu'il faut retenir de ce chapitre~:}
\begin{itemize}
    \item \textbf{Windows et Linux} laissent des artefacts riches.
    \item \textbf{La chronologie} relie les horodatages recalés.
    \item \textbf{On date} chaque étape de l'attaque.
    \item \textbf{Supprimer n'est pas effacer}~: récupération possible.
    \item \textbf{On sépare} le tracé du déduit.
\end{itemize}
\end{boxresume}

\begin{boxinfo}
\textbf{Points de contrôle de l'élève.} Je sais~: citer les artefacts Windows
et Linux~; construire une chronologie~; dater une attaque~; expliquer la
récupération~; séparer trace et interprétation.
\end{boxinfo}

\bigskip

\begin{boxexo}[title={\ding{45}\quad Exercice~12.1. Artefact Windows \niveaubase}]
Quel artefact Windows révèle les programmes lancés au démarrage~?
\end{boxexo}

\begin{boxexo}[title={\ding{45}\quad Exercice~12.2. Artefact Linux \niveaubase}]
Quel artefact Linux révèle les commandes saisies~?
\end{boxexo}

\begin{boxexo}[title={\ding{45}\quad Exercice~12.3. Suppression \niveaubase}]
Pourquoi un fichier supprimé peut-il être récupéré~?
\end{boxexo}

\begin{boxexo}[title={\ding{45}\quad Exercice~12.4. Chronologie \niveaumoyen}]
À partir des faits de FARO, dressez la chronologie des onze jours.
\end{boxexo}

\begin{boxexo}[title={\ding{45}\quad Exercice~12.5. Recaler \niveaumoyen}]
Comment l'écart d'horloge de 40~min modifie-t-il la chronologie~?
\end{boxexo}

\begin{boxexo}[title={\ding{45}\quad Exercice~12.6. Tracé ou déduit \niveaumoyen}]
Classez~: « le service WinUpd32 a été installé le 1\textsuperscript{er}~mars »~;
« l'attaquant cherchait la persistance ».
\end{boxexo}

\begin{boxexo}[title={\ding{45}\quad Exercice~12.7. Persistance \niveaumoyen}]
Quels artefacts trahissent une persistance sous Windows~? Sous Linux~?
\end{boxexo}

\begin{boxexo}[title={\ding{45}\quad Exercice~12.8. Nettoyage \niveaumoyen}]
Pourquoi le « nettoyage » de l'attaquant laisse-t-il souvent des traces~?
\end{boxexo}

\begin{boxexo}[title={\ding{45}\quad Exercice~12.9. Constat système \niveauexpert}]
Rédigez un constat système complet pour l'installation de \texttt{WinUpd32}.
\end{boxexo}

\begin{boxexo}[title={\ding{45}\quad Exercice~12.10. Plan d'analyse \niveauexpert}]
Proposez un plan d'analyse du disque de \texttt{FS-01}, des artefacts à la
chronologie.
\end{boxexo}

\bigskip

\begin{boxinfo}
\textbf{Pour aller plus loin.}
\begin{itemize}
    \item NIST SP~800-86~\cite{nist80086}~; ISO/IEC 27037~\cite{iso27037}.
\end{itemize}
\end{boxinfo}

\begin{boxinfo}
\textbf{Transition.} Le disque ne dit pas tout~: la mémoire vive garde ce qui
n'a jamais été écrit sur disque. Le Chapitre~\ref{chap:memoire} traite du
forensique mémoire.
\end{boxinfo}
